\section{Preliminary results}

We compare the fixed layout with layouts optimized using the oracle model
error, waveform misfit, and illumination objectives. The primary metric is the mean squared model error (MSE) over the complete grid. The multi-seed protocol holds the true and initial models fixed while varying the paired receiver initialization and stochastic FWI source sequence.

% Across ten paired repetitions, every optimized layout reduces final model MSE
% relative to its fixed-layout baseline (Table~\ref{tab:multiseed-summary}). The
% oracle and illumination objectives give mean gains of $18.6\%$ and $16.7\%$,
% respectively, with illumination showing greater variability. Illumination also has the lowest runtime and memory usage, with a mean runtime of \(356.5\,\mathrm{s}\) and a mean peak memory increase of \(9.8\,\mathrm{GiB}\), compared with \(479.2\,\mathrm{s}\) and \(64.3\,\mathrm{GiB}\) for oracle model error and \(503.1\,\mathrm{s}\) and \(66.1\,\mathrm{GiB}\) for waveform misfit. Waveform misfit
% gives a smaller mean gain of $9.4\%$ but the lowest variability. Individual repetitions are
% shown in Appendix Fig.~\ref{fig:paired-improvement}. An illustrative
% reconstruction is shown in Fig.~\ref{fig:reconstruction-comparison}, and the
% corresponding receiver trajectories appear in Appendix
% Fig.~\ref{fig:receiver-trajectories}. The reconstruction using the oracle model error is also shown in Appendix Fig.~\ref{fig:reconstruction-comparison-oracle}.

Across ten paired repetitions, every optimized layout reduces final model MSE
relative to its fixed-layout baseline (Table~\ref{tab:multiseed-summary}).
Oracle and illumination give the largest mean improvements, while illumination
requires substantially less runtime and memory than the objectives
differentiated through FWI. Waveform misfit gives a smaller
but less variable improvement. Figure~\ref{fig:reconstruction-comparison} shows
one representative reconstruction; individual repetitions and receiver
trajectories are reported in the Appendix.


\begin{table}[!htbp]
  \centering
  \scriptsize
  \setlength{\tabcolsep}{3.5pt}
  \caption{Multi-seed reconstruction and runtime summary. Values are
  mean $\pm$ sample standard deviation; improvements are relative to the paired
  fixed-layout run. Peak memory increase is reported as a mean rounded to
  $0.1\,\mathrm{GiB}$. Every reconstruction uses 100 FWI updates.}
  \label{tab:multiseed-summary}
  \begin{tabular}{@{}lrrrrr@{}}
    \toprule
    Method & Final model MSE & Improvement (\%) & Improved & Runtime (s) & Peak $\Delta$Memory (GiB) \\
    \midrule
    Fixed layout & $3358.9\pm65.5$ & $0.0$ & - & $166.3\pm0.9$ & 1.2 \\
    Stochastic illumination & $2796.2\pm274.6$ & $16.7\pm8.6$ & 10/10 & $356.5\pm0.8$ & 9.8 \\
    Waveform misfit & $3041.6\pm125.4$ & $9.4\pm4.3$ & 10/10 & $503.1\pm0.7$ & 66.1 \\
    Oracle model error & $2732.3\pm242.4$ & $18.6\pm7.1$ & 10/10 & $479.2\pm1.2$ & 64.3 \\
    \bottomrule
  \end{tabular}
\end{table}

\begin{figure}[t]
  \centering
  \includegraphics[width=0.85\linewidth]{figures/comparison_plot_full_auto_diff_cropped_colorbar_true_model_outlined.pdf}
  \caption{Reconstruction and receiver-geometry comparison for one
  representative repetition. The dashed cyan box marks the true
  $55\,\mathrm{mm}$ square inclusion; red crosses are the four fixed sources and
  yellow stars are the 30 receivers. The wave-speed scale is shared, and
  percentages give MSE reduction from the fixed layout. Every reconstruction
  uses 100 FWI updates. The color bar is cropped at 80\% of the maximum value of the true model to emphasize differences between the design methods.}
  \label{fig:reconstruction-comparison}
\end{figure}
